
@@M0@@, one has
@@M1@@
where @@M2@@ and @@M3@@ are the structural morphisms of @@M4@@ and @@M5@@. Since @@M6@@, one has:
@@M7@@
by IV 2.12.4.

By (3.4.1), one therefore has:
@@M8@@
Let @@M9@@ and @@M10@@ be the classes of the canonical invertible sheaves on @@M11@@ and @@M12@@. Then, by VI 1.1, @@M13@@ (resp. @@M14@@) has as a basis over @@M15@@ the elements @@M16@@ for @@M17@@ (resp. @@M18@@ for @@M19@@). The projection formula gives:
@@M20@@
If @@M21@@ is the ideal of @@M22@@ in @@M23@@, @@M24@@ is an invertible ideal, and one has @@M25@@. Moreover, @@M26@@ is defined by the ideal @@M27@@ of @@M28@@ that is the kernel of @@M29@@, and is generated by a sub-Module of @@M30@@ isomorphic to @@M31@@; consequently, @@M32@@. The image of @@M33@@ is therefore the @@M34@@-submodule generated by the @@M35@@ for @@M36@@. The homomorphism @@M37@@ is thus an isomorphism, and consequently the same holds for the homomorphism @@M38@@.

\sgaStableTarget{sga6.e09.sec.4}\subsection*{4. Principal bundles under the groups @@M39@@}

\sgaStableTarget{sga6.e09.prop.4-1}\textbf{Proposition 4.1.} \uline{Let} @@M40@@ \uline{be a noetherian scheme}, @@M41@@ \uline{the} @@M42@@-\uline{group} @@M43@@, @@M44@@ \uline{a torsor under} @@M45@@, \uline{and} @@M46@@ \uline{the locally free} @@M47@@\uline{-Module of rank} @@M48@@ \uline{corresponding to} @@M49@@. \uline{Denote by} @@M50@@ \uline{the class of} @@M51@@ \uline{in} @@M52@@, \uline{and by} @@M53@@ \uline{the} @@M54@@-\uline{ideal generated}

